\documentclass{ximera}

\input{../../preamble.tex}

\title{Function Features}


\begin{document}

\begin{abstract}
extreme values
\end{abstract}
\maketitle





Functions are packages containing three sets: domain, codomain, and pairs.  The pairs connect the domain values to codomain values and it is this relationship we want to investigate.   


For us, the codomain is often $\mathbb{R}$ and so we often ignore this until it becomes important.  (There are many reasons for it to become important).  Thus, we often think of our real-valued functions as three sets: domain, range, and pairs. 


We would like to know how the range values are affected by the domain values. 


\begin{itemize}
\item We are interested in the function values - the range values.  But they are controlled by the domain values. 
\item We ask questions about the function values, but the answers are in the domain.
\end{itemize}




\begin{idea} \textbf{\textcolor{blue!55!black}{Extreme Function Values}} 


$\blacktriangleright$ \textbf{\textcolor{purple!85!blue}{Where}} is the maximum value of the function? 

Translation: Which \textbf{\textcolor{purple!85!blue}{domain}} value results in the maximum value of the function? 


$\blacktriangleright$ \textbf{\textcolor{purple!85!blue}{Where}} is the minimum value of the function? 

Translation: Which \textbf{\textcolor{purple!85!blue}{domain}} value results in the minimum value of the function? 


\end{idea}








A graph provides a global view of all of the pairs, which reveals many of the patterns, features, and characteristics we seek.   

However, we must separate the graph from the function.  The graph helps us answer the questions, but the graph doesn't hold the answers.  The answers are in the domain and range.  The answers are not the points in the graph.  The answers are the numbers in the domain and range.



\begin{center}

\textbf{\textcolor{blue!55!black}{The answers are not the points in the graph.}} 


\textbf{\textcolor{blue!55!black}{The answers are the numbers in the domain and range.}}

\end{center}







\pdfOnly{
\textbf{$\mathcal{VIDEO}$: Two Types of Questions}
}
\begin{explanation} \textbf{\textcolor{red!75!green}{Video: Two Types of Questions}}



\begin{onlineOnly}
[ Click on the arrow to the right to expand for the video. ]
\begin{expandable} 

\begin{center}
\youtube{nMK4bo1Ad8g}
\end{center}

\end{expandable}
\end{onlineOnly}

\end{explanation}









The graph displays patterns in its points.  These patterns help us analyze the function which has values in the range connected to numbers in the domain.




\subsection*{Maximum}


\begin{definition} \textbf{\textcolor{green!50!black}{Global Maximum}} 


Let $f$ be a function defined on the domain $D$.

Then $M$ is the \textbf{global maximum value} of $f$ on $D$ if    


\begin{itemize}
\item there exists $d \in D$, such that $f(d) = M$.   i.e., $M$ has to be a function value 
\item $f(d) \leq M$ for all $d \in D$. i.e., all function values are less than or equal to $M$ 

\end{itemize}


If there is no such $M$, then $f$ has no maximum value.


\textbf{Note:} The global maximum is often referred to as just the maximum.

\end{definition}



\begin{observation}\textbf{\textcolor{-red!75}{Graph vs. Function}}

The maximum value is visually encoded in the highest point on the graph of the function.  The second coordinate is the maximum value and the first coordinate is the domain number where the maximum occurs.
\end{observation}










\begin{example} Maximum

Below is the graph of $y=f(x)$.  

\begin{image}
\begin{tikzpicture} 
  \begin{axis}[
            domain=-10:10, ymax=10, xmax=10, ymin=-10, xmin=-10,
            axis lines =center, xlabel=$x$, ylabel={$y=f(x)$},
            grid = both, 
            ytick={-10,-8,-6,-4,-2,2,4,6,8,10}, 
            xtick={-10,-8,-6,-4,-2,2,4,6,8,10},
            yticklabels={$-10$,$-8$,$-6$,$-4$,$-2$,$2$,$4$,$6$,$8$,$10$}, 
            xticklabels={$-10$,$-8$,$-6$,$-4$,$-2$,$2$,$4$,$6$,$8$,$10$},
            ticklabel style={font=\scriptsize},
            every axis y label/.style={at=(current axis.above origin),anchor=south},
            every axis x label/.style={at=(current axis.right of origin),anchor=west},
            axis on top
          ]
          
          \addplot [line width=2, penColor, smooth,samples=100,domain=(-5:0)] ({x},{(x+5)*(1-x)});
          \addplot [line width=2, penColor, smooth,samples=100,domain=(1:5)] ({x},{-x-1});
          \addplot [line width=2, penColor, smooth,samples=100,domain=(5:8)] ({x},{2*x-16});
          

          \addplot [color=penColor,fill=white,only marks,mark=*] coordinates{(0,5) (1,-2) (8,0)};
          \addplot [color=penColor,only marks,mark=*] coordinates{(-5,0) (5,-6)};


           

  \end{axis}
\end{tikzpicture}
\end{image}





\textcolor{gray!90}{alternate text: Cartesian plane. Horizontal axis labeled x. Vertical axis labeled y=f(x).  The graph consists of two pieces. The first piece begins with a solid dot at (-5,0), moves up to (-2,9), then moves down to a hollow dot at (0,5). The second piece has a 'Vee' shape. It begins with a hollow dot at (1,-2), moves down to the corner at (5,-6), then up to a hollow dot at (8,0).}

\pdfOnly{
\textbf{$\mathcal{DESMOS}$: Maximum Value}
}
\begin{onlineOnly}
\textbf{\textcolor{blue!55!black}{$\blacktriangleright$ desmos graph}} 
\begin{center}
\desmos{dsxpdq3ooo}{400}{300}
\end{center}
\end{onlineOnly}



\begin{question}

The highest point on the graph of $f$ is


\begin{multipleChoice}
\choice {$(-5,0)$}
\choice [correct]{$(-2,9)$}
\choice {$(0,5)$}
\choice {$(-1,-2)$}
\choice {$(5,-6)$}
\choice {$(8,0)$}
\end{multipleChoice}
\end{question}


\begin{question}


The maximum value of $f$ is $\answer{9}$, which occurs at $\answer{-2}$.
\end{question}







\end{example}
















\newpage



\begin{example} Maximum

Below is the graph of $y=g(t)$.  

\begin{image}
\begin{tikzpicture} 
  \begin{axis}[
            domain=-10:10, ymax=10, xmax=10, ymin=-10, xmin=-10,
            axis lines =center, xlabel=$t$, ylabel={$y=g(t)$},
            grid = both, 
            ytick={-10,-8,-6,-4,-2,2,4,6,8,10}, 
            xtick={-10,-8,-6,-4,-2,2,4,6,8,10},
            yticklabels={$-10$,$-8$,$-6$,$-4$,$-2$,$2$,$4$,$6$,$8$,$10$}, 
            xticklabels={$-10$,$-8$,$-6$,$-4$,$-2$,$2$,$4$,$6$,$8$,$10$},
            ticklabel style={font=\scriptsize},
            every axis y label/.style={at=(current axis.above origin),anchor=south},
            every axis x label/.style={at=(current axis.right of origin),anchor=west},
            axis on top
          ]
          

          \addplot [line width=2, penColor, smooth,samples=100,domain=(-6:5)] ({x},{-x-1});
          \addplot [line width=2, penColor, smooth,samples=100,domain=(5:8)] ({x},{2*x-16});

          \addplot [color=penColor,fill=white,only marks,mark=*] coordinates{(-6,5)};
          \addplot [color=penColor,only marks,mark=*] coordinates{(5,-6) (8,0)};


           

  \end{axis}
\end{tikzpicture}
\end{image}





\textcolor{gray!90}{alternate text: Cartesian plane. Horizontal axis labeled t. Vertical axis labeled y=g(t).  The graph has a 'Vee' shape. It begins with a hollow dot at (-6,5), moves down to the corner at (5,-6), then up to a hollow dot at (8,0).}

\pdfOnly{
\textbf{$\mathcal{DESMOS}$: Maximum Value}
}
\begin{onlineOnly}
\textbf{\textcolor{blue!55!black}{$\blacktriangleright$ desmos graph}} 
\begin{center}
\desmos{enrjrmglxn}{400}{300}
\end{center}
\end{onlineOnly}




There is no highest point on the graph.  Therefore, $g$ has no maximum value.

\end{example}


































\subsection*{Minimum}


\begin{definition} \textbf{\textcolor{green!50!black}{Global Minimum}} 


Let $f$ be a function defined on the domain $D$.

Then $M$ is the \textbf{global minimum value} of $f$ on $D$ if    


\begin{itemize}
\item there exists $d \in D$, such that $f(d) = M$.   i.e., $M$ has to be a function value 
\item $f(d) \geq M$ for all $d \in D$. i.e., all function values are greater than or equal to $M$. 

\end{itemize}


If there is no such $M$, then $f$ has no minimum value.



\textbf{Note:} The global minimum is often referred to as just the minimum.

\end{definition}



\begin{observation} \textbf{\textcolor{-red!75}{Graph vs. Function}}

The minimum value is visually encoded in the lowest point on the graph of the function.  The second coordinate is the minimum value and the first coordinate is the domain number where the minimum occurs.
\end{observation}






\newpage




\begin{example} Minimum

Below is the graph of $y=f(x)$.  

\begin{image}
\begin{tikzpicture} 
  \begin{axis}[
            domain=-10:10, ymax=10, xmax=10, ymin=-10, xmin=-10,
            axis lines =center, xlabel=$x$, ylabel={$y=f(x)$},
            grid = both, 
            ytick={-10,-8,-6,-4,-2,2,4,6,8,10}, 
            xtick={-10,-8,-6,-4,-2,2,4,6,8,10},
            yticklabels={$-10$,$-8$,$-6$,$-4$,$-2$,$2$,$4$,$6$,$8$,$10$}, 
            xticklabels={$-10$,$-8$,$-6$,$-4$,$-2$,$2$,$4$,$6$,$8$,$10$},
            ticklabel style={font=\scriptsize},
            every axis y label/.style={at=(current axis.above origin),anchor=south},
            every axis x label/.style={at=(current axis.right of origin),anchor=west},
            axis on top
          ]
          
          \addplot [line width=2, penColor, smooth,samples=100,domain=(-5:0)] ({x},{(x+5)*(1-x)});
          \addplot [line width=2, penColor, smooth,samples=100,domain=(1:5)] ({x},{-x-1});
          \addplot [line width=2, penColor, smooth,samples=100,domain=(5:8)] ({x},{2*x-16});

          \addplot [color=penColor,fill=white,only marks,mark=*] coordinates{(0,5) (1,-2) (8,0)};
          \addplot [color=penColor,only marks,mark=*] coordinates{(-5,0) (5,-6)};


           

  \end{axis}
\end{tikzpicture}
\end{image}






\textcolor{gray!90}{alternate text: Cartesian plane. Horizontal axis labeled x. Vertical axis labeled y=f(x).  The graph consists of two pieces. The first piece begins with a solid dot at (-5,0), moves up to (-2,9), then moves down to a hollow dot at (0,5). The second piece has a 'Vee' shape. It begins with a hollow dot at (1,-2), moves down to the corner at (5,-6), then up to a hollow dot at (8,0).}

\pdfOnly{
\textbf{$\mathcal{DESMOS}$: Minimum Value}
}
\begin{onlineOnly}
\textbf{\textcolor{blue!55!black}{$\blacktriangleright$ desmos graph}} 
\begin{center}
\desmos{dsxpdq3ooo}{400}{300}
\end{center}
\end{onlineOnly}




\begin{question} 


The lowest point on the graph of $f$ is


\begin{multipleChoice}
\choice {$(-5,0)$}
\choice {$(-2,9)$}
\choice {$(0,5)$}
\choice {$(-1,-2)$}
\choice [correct]{$(5,-6)$}
\choice {$(8,0)$}
\end{multipleChoice}
\end{question}


\begin{question} 


The minimum value of $f$ is $\answer{-6}$, which occurs at $\answer{5}$.
\end{question}







\end{example}






\newpage




\begin{example} Minimum

Below is the graph of $y=H(k)$.  

\begin{image}
\begin{tikzpicture} 
  \begin{axis}[
            domain=-10:10, ymax=10, xmax=10, ymin=-10, xmin=-10,
            axis lines =center, xlabel=$k$, ylabel={$y=H(k)$},
            grid = both, 
            ytick={-10,-8,-6,-4,-2,2,4,6,8,10}, 
            xtick={-10,-8,-6,-4,-2,2,4,6,8,10},
            yticklabels={$-10$,$-8$,$-6$,$-4$,$-2$,$2$,$4$,$6$,$8$,$10$}, 
            xticklabels={$-10$,$-8$,$-6$,$-4$,$-2$,$2$,$4$,$6$,$8$,$10$},
            ticklabel style={font=\scriptsize},
            every axis y label/.style={at=(current axis.above origin),anchor=south},
            every axis x label/.style={at=(current axis.right of origin),anchor=west},
            axis on top
          ]
          

          \addplot [line width=2, penColor, smooth,samples=100,domain=(-6:5)] ({x},{-x-1});
          \addplot [line width=2, penColor, smooth,samples=100,domain=(5:8)] ({x},{2*x-16});

          \addplot [color=penColor,fill=white,only marks,mark=*] coordinates{(-6,5) (5,-6)};
          \addplot [color=penColor,only marks,mark=*] coordinates{(8,0)};


           

  \end{axis}
\end{tikzpicture}
\end{image}






\textcolor{gray!90}{alternate text: Cartesian plane. Horizontal axis labeled k. Vertical axis labeled y=H(k).  The graph has a 'Vee' shape. It begins with a hollow dot at (-6,5), moves down to the corner at (5,-6), then up to a hollow dot at (8,0).}

\pdfOnly{
\textbf{$\mathcal{DESMOS}$: Maximum Value}
}
\begin{onlineOnly}
\textbf{\textcolor{blue!55!black}{$\blacktriangleright$ desmos graph}} 
\begin{center}
\desmos{adhfbavgxd}{400}{300}
\end{center}
\end{onlineOnly}



There is no lowest point on the graph.  Therefore, $H$ has no minimum value.

\end{example}



\subsection*{Where?}


\textbf{\textcolor{blue!55!black}{$\blacktriangleright$}} A function might not have a global maximum value.

If a function has a global maximum value, it has exactly one global maximum value.

However, this one global maximum value might occur at many domain numbers.



\textbf{\textcolor{blue!55!black}{$\blacktriangleright$}} A function might not have a global minimum value.

If a function has a global minimum value, it has exactly one global minimum value.

However, this one global minimum value might occur at many domain numbers.











\begin{example} A Weird Example

Below is the graph of $y=C(x)$.  

\begin{image}
\begin{tikzpicture} 
  \begin{axis}[
            domain=-10:10, ymax=10, xmax=10, ymin=-10, xmin=-10,
            axis lines =center, xlabel=$x$, ylabel={y},
            grid = both, 
            ytick={-10,-8,-6,-4,-2,2,4,6,8,10}, 
            xtick={-10,-8,-6,-4,-2,2,4,6,8,10},
            yticklabels={$-10$,$-8$,$-6$,$-4$,$-2$,$2$,$4$,$6$,$8$,$10$}, 
            xticklabels={$-10$,$-8$,$-6$,$-4$,$-2$,$2$,$4$,$6$,$8$,$10$},
            ticklabel style={font=\scriptsize},
            every axis y label/.style={at=(current axis.above origin),anchor=south},
            every axis x label/.style={at=(current axis.right of origin),anchor=west},
            axis on top
          ]
          

          \addplot [line width=2, penColor, smooth,samples=100,domain=(-6:5)] ({x},{2});

          \addplot [color=penColor,only marks,mark=*] coordinates{(-6,2)};
          \addplot [color=penColor,only marks,mark=*] coordinates{(5,2)};


           

  \end{axis}
\end{tikzpicture}
\end{image}


\textcolor{gray!90}{alternate text: Cartesian plane. Horizontal axis labeled x. Vertical axis labeled y.  The graph is a horizontal line segment connecting solid dots at (-6,2) and (5,2).}




$C$ is a constant function with $[-6,5]$ as its domain.

$C$ has only one function value.  That is $2$.


\textbf{\textcolor{blue!55!black}{$\blacktriangleright$}} $C(x) \leq 2$ for every domain number.

So, $2$ is the global maximum value of $C$, and it occurs at every domain number.


\textbf{\textcolor{blue!55!black}{$\blacktriangleright$}} $C(x) \geq 2$ for every domain number.

So, $2$ is the global minimum value of $C$, and it occurs at every domain number.



\end{example}


Functions can do weird things.










\begin{onlineOnly}
\begin{center}
\textbf{\textcolor{green!50!black}{ooooo-=-=-=-ooOoo-=-=-=-ooooo}} \\

more examples can be found by following this link\\ \link[More Examples of Function Graphs]{https://ximera.osu.edu/csccmathematics/precalculus/precalculus/functionGraphs/examples/exampleList}

\end{center}
\end{onlineOnly}




\end{document}
